\documentclass[11pt, a4paper]{article}

\usepackage[T1]{fontenc}
\usepackage[utf8]{inputenc}
\usepackage{tgpagella}
\usepackage[spanish, es-noshorthands]{babel}
\usepackage{amsmath, amssymb, bm}
\usepackage{tabularx}
\usepackage[table, xcdraw]{xcolor}
\usepackage[most]{tcolorbox}
\usepackage[margin=2.5cm]{geometry}
\usepackage{fancyhdr}

\pagestyle{fancy}
\fancyhf{}
\setlength{\headheight}{16pt}
\lhead{\textbf{UCB Sede Tarija} | Ing. Mecatrónica}
\chead{}
\rhead{IMT-121 Circuitos Electrónicos I | Instructor Key}
\lfoot{Prof: M.Sc. Ing. Bernardo Quiroga Turdera}
\rfoot{Página \thepage}
\renewcommand{\headrulewidth}{0.4pt}

\definecolor{ucbblue}{RGB}{0, 51, 102}

\begin{document}

\begin{tcolorbox}[colback=ucbblue!5!white, colframe=ucbblue, title=\textbf{Clave del Instructor: Hacia la Superposición}]
    \centering \textbf{Sesión Asincrónica (Semana 5) / Revisión Presencial (Semana 6)}
\end{tcolorbox}

\section*{Solucionario y Guía de Intervención (Actividad Formativa)}

\subsection*{Actividad 1}
\textbf{Respuesta correcta:} a) La red física (conductancias y topología). b) Excitaciones externas (corrientes inyectadas). c) No. Si solo cambia la fuente, la estructura y resistencias se mantienen, dejando $\bm{A}$ fija. d) Predicción: $[4, 2]^T\,\text{V}$. \\
\textbf{Error frecuente:} "Se duplica porque 10 es el doble de 5." \\
\textbf{Prompt correctivo:} "¿Qué condición física del circuito garantiza que esa regla matemática se cumpla siempre?" (R: Mismo modelo lineal fijo).

\subsection*{Actividad 2 (Homogeneidad)}
\begin{itemize}
    \item $k=2 \implies [4, 2]^T\,\text{V}$
    \item $k=-1 \implies [-2, -1]^T\,\text{V}$ (Invierte referencias de signo)
    \item $k=0.5 \implies [1, 0.5]^T\,\text{V}$
\end{itemize}
\textbf{Explicación esperada:} En una red lineal cuya matriz $\bm{A}$ permanece fija, escalar la excitación por $k$ multiplica todas las respuestas por $k$. \\
\textbf{Error frecuente:} Omitir la condición de "red fija". 

\subsection*{Actividad 3 (Additividad)}
\textbf{Suma excitaciones:} $\bm{b}_1 + \bm{b}_2 = [5, 3]^T\,\text{mA}$. \\
\textbf{Predicción de respuesta:} $\bm{x}_1 + \bm{x}_2 = [2.6, 2.8]^T\,\text{V}$. \\
\textbf{Propiedad observada:} Additividad. \\
\textbf{Explicación conceptualmente débil:} "Se suman porque las matemáticas lo dicen". \\
\textbf{Prompt correctivo:} "Ambas respuestas derivan de multiplicar por la inversa de la misma matriz constante. Físicamente, estamos superponiendo efectos sobre el mismo sistema inalterado".

\subsection*{Actividad 4 (Clasificación)}
\textbf{Respuestas:} 
1: \textbf{H} \quad 
2: \textbf{A} \quad 
3: \textbf{L} \quad 
4: \textbf{?} (Al cambiar resistor, $\bm{A}$ cambia; se pierde la referencia lineal directa).

\subsection*{Actividad 5 (Puente a la Sesión Presencial de Semana 6)}
\textbf{Hipótesis esperada:} La respuesta total debería ser la suma algebraica de las contribuciones individuales que genera cada fuente. \\
\textbf{Lo que falta aprender (Target Cognitivo):} "Todavía no sabemos qué hacer físicamente con el resto de las fuentes en el circuito cuando intentamos simular que 'solo hay una fuente activa'". \\
\textit{Nota para el Instructor: No exigir ni revelar aquí las reglas de desactivación (cortocircuito / circuito abierto). Esta pregunta está diseñada puramente para generar la necesidad del Teorema de Superposición en la Sesión 1 de la Semana 6.}

\vspace{0.5cm}
\begin{tcolorbox}[colback=white, colframe=black, boxrule=0.5pt, arc=1mm]
\textbf{Prompt de Transición en Aula (10 min iniciales, Semana 6, S01):}
\begin{enumerate}
    \item (Proyectar $\bm{A}\bm{x} = \bm{b}$) "Si $\bm{A}$ es fija y $\bm{b}$ cambia a $-2\bm{b}$, ¿qué ocurre con $\bm{x}$?" (Espera: Homogeneidad $\rightarrow -2\bm{x}$).
    \item "Dado $\bm{x}_1=[2,1]^T$ y $\bm{x}_2=[0.6,1.8]^T$. Sin resolver sistema, respondan a excitación $2\bm{b}_1 - \bm{b}_2$". (Espera: Linealidad combinada $\rightarrow [3.4, 0.2]^T$).
    \item \textbf{Cierre:} "Si podemos separar y sumar contribuciones matemáticas, hagámoslo en el esquemático físico. ¿Cómo apagamos una fuente de voltaje sin romper las Leyes de Kirchhoff? (Inicia lección Superposición)".
\end{enumerate}
\end{tcolorbox}

\end{document}
